%! AP Statistics — Unit 1, Lesson 1.2 — Group Activity KEY
\documentclass[10pt]{article}
\usepackage{apstats-article}
\usepackage{apstats-key}

%=============================================================================
\begin{document}

\pageheader{Unit 1, Lesson 1.2}{Group Activity --- Answer Key}
\begin{tabularx}{\linewidth}{X X X}
Name:\;\ans{Key} & Partner:\; & Period:\;
\end{tabularx}

\vspace{0.14in}

% ── PART 1 ───────────────────────────────────────────────────────────────────
{\large\bfseries\color{navy} Part 1 \quad Mini-Whiteboard Round}

\vspace{0.1in}

\renewcommand{\arraystretch}{2.0}
\begin{tabularx}{\linewidth}{@{} c >{\raggedright\arraybackslash}p{0.36\linewidth}
                                    C{0.18\linewidth} X @{}}
\rowcolor{sky}
\textbf{\#} & \textbf{Variable} & \textbf{Answer} & \textbf{Reason} \\
\midrule
1 & Favorite season
  & \ans{C}  & \ans{Labels a category; averaging seasons is meaningless} \\
2 & Push-ups in one minute
  & \ans{QD} & \ans{Count of whole-number reps; gaps between values} \\
3 & Blood type
  & \ans{C}  & \ans{Places individuals into distinct categories (A, B, AB, O)} \\
4 & Distance driven to school
  & \ans{QC} & \ans{Measured distance; can be any non-negative real number} \\
5 & Number of Instagram followers
  & \ans{QD} & \ans{Count of whole-number followers; gaps between values} \\
6 & Jersey number
  & \ans{C}  & \ans{Numeric label; averaging jersey numbers is meaningless} \\
\end{tabularx}

\vspace{0.18in}

% ── PART 2 ───────────────────────────────────────────────────────────────────
{\large\bfseries\color{navy} Part 2 \quad Card Sort}

\vspace{0.1in}

\renewcommand{\arraystretch}{2.2}
\begin{tabularx}{\linewidth}{@{} |>{\centering\arraybackslash}X
                                    |>{\centering\arraybackslash}X
                                    |>{\centering\arraybackslash}X| @{}}
\hline
\rowcolor{greenbg}
\textbf{\textcolor{greenacc}{Categorical}} &
\textbf{\textcolor{navy}{Quantitative Discrete}} &
\textbf{\textcolor{redacc}{Quantitative Continuous}} \\
\hline
\ans{A — Eye color} &
\ans{B — Books read} &
\ans{C — Weight of watermelon} \\
\ans{D — ZIP code} &
\ans{E — Pets owned} &
\ans{F — Annual salary} \\
\ans{G — Streaming service} &
\ans{K — Siblings} &
\ans{H — Hours of sleep} \\
\ans{L — Country of citizenship} &
 &
\ans{J — Height (cm)} \\
\ans{I — Star rating (debatable)} &  & \\
\hline
\end{tabularx}

\vspace{0.1in}

\begin{teachernote}
\textbf{Card I (star rating):} Accept Categorical or QD if the student provides reasoning.
Most AP Statistics resources treat ordinal ratings (1--5 stars) as categorical because
equal spacing between levels cannot be assumed. However, many analysts treat them as
quantitative for convenience. The key teaching point is that \emph{context and judgment
matter} --- the classification is not always automatic.\\[4pt]
\textbf{Card H (hours of sleep):} Some students will say discrete because they think of
sleep in whole hours. Clarify that sleep is a measured duration --- it can be 6.5 hours,
7.25 hours, etc. --- so it is continuous.\\[4pt]
\textbf{Card F (annual salary):} Technically discrete (paid in whole cents), but so many
possible values that it is treated as continuous in practice. Accept either with justification.
\end{teachernote}

\vspace{0.1in}

% ── PART 3 ───────────────────────────────────────────────────────────────────
{\large\bfseries\color{navy} Part 3 \quad Discuss \& Decide}

\vspace{0.1in}

\begin{tcolorbox}[colback=redbg, colframe=redacc, boxrule=0.8pt, arc=2mm,
    left=3mm, right=3mm, top=2mm, bottom=2mm]
\small\textbf{Contested variable: Rating of a movie (1--5 stars) --- Card I}

\begin{enumerate}[leftmargin=1.5em, itemsep=6pt]
  \item Can you compute the average? Does it mean something?\\[4pt]
        \ans{Yes, you can compute it arithmetically. Whether it is meaningful is debatable:
        a 3.7-star average conveys something useful in everyday language, but the difference
        between 1 and 2 stars may not equal the difference between 4 and 5 stars.}

  \item Are the gaps between ratings equal?\\[4pt]
        \ans{Not necessarily. A viewer who rates a film 1 star may feel far more strongly
        negative than the gap to 2 stars implies. Equal spacing is an assumption, not a fact.
        This is what makes ordinal scales ambiguous.}

  \item Final classification: \ans{Categorical (ordinal)} or \ans{QD (with assumptions)}\\
        Reason: \ans{Accept either; look for awareness that the choice depends on whether equal spacing is assumed.}
\end{enumerate}
\end{tcolorbox}

\vspace{0.14in}

\begin{tcolorbox}[colback=goldbg, colframe=goldacc, boxrule=0.8pt, arc=2mm,
    left=3mm, right=3mm, top=2mm, bottom=2mm]
\small\textbf{Reflection:}\\[6pt]
\ans{Most groups contest Card I (star rating) or Card H (sleep). Resolution usually comes from
asking ``Can this variable take a value between two consecutive values?'' and ``Are the gaps
between levels equal?'' Accept any response that shows the group reasoned through it rather
than guessing.}
\end{tcolorbox}

\end{document}
